\section{Почему \TeX{} всё ещё важен в 2026 году}

Когда Дональд Кнут выпустил первую версию \TeX{} в 1978 году, у него была одна цель:
красиво сверстать собственные книги. \citetext{knuth1984tex}{Спустя сорок восемь лет
его творение остаётся золотым стандартом математической и научной вёрстки.}

\subsection{Проблема точности}

Текстовые процессоры и браузеры оптимизированы под перекомпоновку, а не под точность.
Абзац в Google Docs выглядит по-разному на каждом экране.
\citetext{bringhurst2004elements}{\TeX{}, напротив, решает задачу глобальной
оптимизации: он вычисляет разрывы строк для целого абзаца сразу, достигая
результатов, которые профессиональные типографы описывают как ``невозможно получить
вручную''.}

Разница измерима. \citetext{ruder2013typography}{В исследовании 2013 года,
сравнивавшем скорость чтения на разных движках вёрстки, текст, набранный в
\TeX{}, показал превосходство на 12\% в задачах длительного чтения.}

\subsection{Фактор формул}

Для тех, кто пишет математику, \TeX{} --- не выбор, а инфраструктура. Встроенные
формулы вроде $E = mc^2$ и выносные выражения

\[
\int_0^1 x^2 \, dx = \frac{1}{3}
\]

\citetext{lamport1994latex}{Ни одна другая система не обеспечивает такой
согласованности между форматами вывода: бумага, PDF и браузеры с поддержкой MathML.}

\subsection{Статические сайты и ренессанс \TeX{}}

Pandoc, универсальный конвертер документов, сделал практичным авторство в \TeX{} и
публикацию в вебе. \citetext{pandoc2026}{Инструмент переводит исходники \LaTeX{} в
HTML5 с MathML, сохраняя и структуру, и семантику.} В сочетании с генератором
статических сайтов это даёт авторам лучшее из двух миров: качество вёрстки \TeX{}
на этапе написания и охват веба на этапе публикации.

\begin{center}
\includegraphics[width=0.82\linewidth]{assets/random-grid.png}
\end{center}

\subsection{Зачем всё это?}

Скептики спрашивают: если уже есть Medium и Notion, зачем писать в \TeX{}?
\citetext{knuth1984tex}{Ответ --- долговечность. Файл \texttt{.tex}, написанный в
1984 году, компилируется сегодня и будет компилироваться в 2034 году. Проприетарные
форматы приходят и уходят; plain-text исходники с явной разметкой выживают.}

Помимо долговечности, есть ремесло. Писать в \TeX{} значит думать о структуре
раньше, чем о стиле, --- дисциплина, которая улучшает и прозу, и подачу.

\subsection{Что дальше}

\citetext{pandoc2026}{Экосистема продолжает расти. Современные движки вроде
\texttt{xelatex} поддерживают Unicode и OpenType-шрифты нативно, инструменты
\texttt{latexmk} автоматизируют многопроходную компиляцию, а проекты вроде этого
сайта демонстрируют, что \TeX{} и веб --- не конкуренты, а взаимодополняющие
технологии.}

\begin{itemize}
  \item \TeX{} --- точная вёрстка и математика
  \item Pandoc --- универсальная конвертация форматов
  \item Статические генераторы --- быстрая, безопасная, версионируемая публикация
  \item MathML --- доступная математика в вебе
\end{itemize}

Это стек, который я выбрал для этого сайта. Он работает.
